\section{Limitations and conclusion}
\label{sec:limitations}

Seven limitations should frame any use of this model, and each is disclosed in the codebase as well as here.

\paragraph{No scenario result has been replicated.} This is the limitation that matters most, because scenarios are what the model is for. No machine-readable numeric scenario results are published for version~1.1, so targets 2 and 3 are closed at the achievable ceiling rather than passed (Section~\ref{sec:targets23}). What is gated is scenario \emph{design} --- the exact policy-switch set each cached scenario toggles --- plus two coarse anchors from a pre-1.0 conference version that are regression pins, not agreement tests, one of which admits values contradicting the sentence it is named for. Nothing in this paper establishes that DEFINE-UK's scenario \emph{outputs} reproduce anything published. The condition for lifting this is specific: the authors publish scenario tables, or either SSRN paper becomes accessible.

\paragraph{The baseline is not a forecast, and its levels are far from outturns.} 2025 real GDP growth of 4.66 per cent against an ONS 1.31; unemployment below the outturn; 2024 emissions some 30 MtCO$_2$e above the DESNZ actuals. The manual itself says the baseline ``should also not be seen as a prediction or forecast''. The adapter therefore serves deltas and never levels, and no DEFINE-UK number appears in any of the suite's forecast surfaces. Users who want UK levels should use the suite's OBR emulator or its Bank of England SVAR.

\paragraph{One published match is a tolerance, not an agreement.} Target 1a passes, and the 0.06\,pp matches at 2030 and 2040 are real. The 2025 growth pass is not: it clears a $\pm0.31$\,pp gate set after the run by 0.01\,pp, on an annualisation convention Table~4 does not state precisely enough to resolve. It is reported as a pass of the gate and not as a match.

\paragraph{The emissions path diverges, and the cause is not established.} The pinned commit runs 3.5 per cent below the published table in 2025 and 23.3 per cent below by 2040. Two checks rule out extraction error on this side, which leaves a vintage or calibration gap between the published table and this commit --- but which vintage produced Table~4 is not known here, and the question is with the authors. Anyone using the adapter's emissions deltas should know that the level path they are differenced from is not the published one.

\paragraph{The multiplier is ours and is not comparable.} 1.78 is a quantity this replication defined. Its horizon and closure place it above any published short-horizon, impulse-based multiplier; its denominator mixes the impulse with endogenous items, and the direction of that bias is not established (Section~\ref{sec:multiplier}). An impulse-based number has not been computed, and choosing the impulse definition is the obvious next piece of work. A previously quoted IMF comparator has been withdrawn as unlocatable.

\paragraph{Baseline and policy runs use different $\alpha_{PP}$.} The upstream Baseline runs production's own-input coefficient \texttt{alpha\_PPLR} at 0.2748; every policy run uses 0.2762 (printed in each run's \texttt{Parameters-sorted\_*.csv}). No policy switch acts before 2027Q1, so the 2023--26 deltas are this parameter change alone, not policy effects; real GDP sits about $-0.16$\,\% of baseline in 2025--26. After 2026 the two cannot be separated, so the published numbers stand and the range below bounds the policy-only effect (the parameter effect stopping in 2027, or persisting at its 2026 level or share of baseline): multiplier 1.78 $\rightarrow$ 1.73--1.75; green-investment GDP peak $+0.92$\,\% $\rightarrow$ $+0.92$--1.08\,\%; housing-regulation real GDP in 2035 $-2.3$ $\rightarrow$ $-0.96$ to $-2.3$; fossil-fuel-ban real GDP in 2030 $-3.85$ $\rightarrow$ $-2.65$ to $-3.85$; emissions deltas change by under 0.25\,\%. Sized in define-uk-model \#20; the fix is not yet decided.

\paragraph{The reimplementation is early, and the model is local-only until it is not.} Milestone~1 is passed and milestone~2 is in progress with 118 of the manual's equations transcribed and its oracle gate unrun; \S3.3.3, the seven sectoral sections, the ecosystem block and the scenarios remain. Until it is complete, DEFINE-UK cannot be hosted at all, because the upstream is unlicensed. Two operational caveats attach: the full-notebook runtime and the upstream's R dependency set are untested in continuous integration, and the oracle half of the validation suite runs only where a cached pinned run exists.

\paragraph{Not an official product.} This is an independent open-source replication and reimplementation, built from the published manual and papers. It is not produced by, endorsed by, or speaking for the DEFINE team; where this work and the publications disagree, the publications are ground truth and the disagreement is recorded as such. The model design belongs to its authors; the implementation is ours.

\subsection*{Conclusion}

The useful thing about DEFINE-UK for a replication programme is that it fails in an instructive way. The code runs, unmodified, at a pinned commit. One published table's macro block reproduces, two of its rows to 0.06\,pp. Another row of the same table diverges by 23 per cent at the horizon. And the results that would matter most --- what happens to output, emissions, employment and balance sheets under a fossil-fuel ban or a green investment programme --- cannot be checked at all, not because the model is wrong but because the numbers were never published.

That combination is common and it is usually reported as a pass. The alternative this paper argues for is a record with three separations held throughout: between checks that can fail and checks that cannot, at the point of claim rather than in a footnote; between a number we computed and a number somebody published; and between a ceiling imposed by the publications and a limitation of the replication, with the lifting condition named in each case. Applied here, that record supports a narrow set of claims --- a baseline macro block reproduced, an emissions divergence pinned, scenario design verified, deltas served with mandatory caveats, and a clean-room reimplementation 118 equations into its second milestone --- and refuses a much larger set that the same underlying work could have been made to look like.

Thirteen defects in a published manual, four of them re-characterised and one claim withdrawn on audit, is the other half of the return. Reimplementing published equations is expensive, and its usual justification is licensing. The justification on display here is different: it is the only reading of a model close enough to find out that a section as published cannot be simulated forward, that a price identity is off by a factor of three because an equation is missing from the document, and that one table's capital block is deflated at the wrong price. None of that was visible from running the code.
